\chapter{Literature Review}

\section{Relationship with Undergraduate Studies}

This internship project synthesizes core theoretical and practical principles from the undergraduate Computer Science \& Engineering curriculum at Independent University, Bangladesh (IUB). The engineering tasks performed during the 3-month placement at Radiant Pharmaceuticals Limited map directly to key academic course domains:

\begin{itemize}
    \item \textbf{Database Management Systems (CSC303):} Designing relational star-schema tables, enforcing foreign key integrity, creating composite indexes, and writing optimized SQL aggregation queries.
    \item \textbf{Object-Oriented Programming (CSC203):} Structuring Django models, custom authentication backends, and view controllers using clean Python object-oriented design patterns.
    \item \textbf{Web Engineering (CSC343):} Developing REST API endpoints, implementing responsive web interfaces, and managing client-server HTTP request-response cycles.
    \item \textbf{Software Engineering (CSC305):} Applying Work Breakdown Structures (WBS), Gantt chart scheduling, requirement analysis, automated unit testing, and software maintainability practices.
\end{itemize}

\section{Related Works}

Academic and industry literature highlights various strategies for enterprise software development, relational query optimization, and e-commerce platform architecture:

\subsection{Web-Based E-Commerce and Enterprise Architecture}
Junhua~\cite{junhua2011design} investigated the design and implementation of web-based e-commerce systems, emphasizing that modular multi-tier architectures provide the necessary separation between presentation, application logic, and backend database storage. Implementing a clean separation of concerns reduces software defect rates and facilitates future system scalability, a structural principle adopted in the dual-application partitioning of this automotive platform.

\subsection{Relational Database Architecture and Query Optimization}
Vicente et al.~\cite{vicente2021rdbms} evaluated database architectures and query performance in web applications, demonstrating that while Object-Relational Mapping (ORM) tools simplify standard CRUD transactions, complex analytical aggregations over large relational datasets suffer measurable latency due to object hydration and redundant join overhead. Bypassing standard ORM layers using direct, indexed SQL query execution minimizes memory consumption and delivers sub-20ms query response times, a strategy implemented in the executive reporting engine of this project.

\subsection{E-Commerce Development and Usability Bottlenecks}
Ullah et al.~\cite{ullah2016developing} analyzed practical methodologies for developing e-commerce web applications, highlighting the importance of responsive catalog navigation, robust session management, and secure transactional workflows for customer retention. Abdullah et al.~\cite{abdullah2021evaluating} evaluated developer usability and execution overhead across e-commerce content management systems, noting that off-the-shelf monolithic CMS platforms introduce significant runtime friction and unnecessary database complexity. In this project, custom lightweight Django templates and REST endpoints are engineered to avoid CMS bloat while maintaining responsive user experience.

\subsection{Recommendation Engines and Search Filtering Performance}
Savadatti and Krishna~\cite{savadatti2024implementation} reviewed collaborative filtering and algorithmic recommendation systems in e-commerce, noting that machine-learning recommendation models frequently introduce cold-start latency and computational bottlenecks when evaluating newly cataloged inventory items. To eliminate cold-start delays and deliver instant vehicle discovery, the platform incorporates an indexed database search and attribute filter engine that provides real-time catalog exploration without recommendation overhead.
